% Source: https://artofproblemsolving.com/community/c3h2928477_stop_the_boredom
% Stable identity: aops-topic-2928477-q09
At most one queen can occupy any row, so there can be at most eight. Eight are attainable: place queens in columns $1,5,8,6,3,7,2,4$ of rows $1$ through $8$, respectively. Columns are distinct. The row-plus-column values are $2,7,11,10,8,13,9,12$, all distinct, and the row-minus-column values are $0,-3,-5,-2,2,-1,5,4$, also all distinct. Thus no diagonal contains two queens. The maximum is $\boxed8$. 
\begin{asy}
usepackage("amsmath");
import geometry;
size(260);
int[] cols={1,5,8,6,3,7,2,4};for(int i=0;i<8;++i)for(int j=0;j<8;++j)filldraw((j,i)--(j+1,i)--(j+1,i+1)--(j,i+1)--cycle,((i+j)%2==0?lightgray:white),gray);for(int i=0;i<8;++i){label("Q",(cols[i]-.5,7.5-i));label(string(i+1),(-.3,7.5-i));label(string(i+1),(i+.5,8.3));}
\end{asy}
